\begin{abstract}
Selecting the LoRA rank for diffusion model fine-tuning is a practical systems problem: higher ranks increase trainable capacity and adaptation cost, while lower ranks improve efficiency but may underfit. We present an empirical engineering study of LoRA rank trade-offs under controlled training budgets in diffusion fine-tuning. Our primary evaluation uses a DDPM U-Net backbone on CIFAR-10 with LoRA ranks \{2, 4, 8, 16, 32\} under fixed optimization settings, and reports parameter counts, runtime, GPU memory, and image quality measured by FID. We then strengthen the analysis with two targeted validations: (i) an extended-budget DDPM setting (ranks 4, 8, 16; 20 epochs), and (ii) a lightweight Tiny DiT validation (ranks 4, 8, 16; 10 epochs). All experiments use a local-folder \texttt{pytorch-fid} protocol with a reproducible CIFAR-10 test reference and 2000 generated evaluation samples per rank.

Across controlled DDPM runs, moderate ranks provide the most favorable efficiency-quality trade-off: rank 4 attains the best FID (\num{124.1380}), rank 8 is competitive (\num{124.2136}), and higher ranks show diminishing returns despite larger trainable parameter budgets. In extended-budget DDPM validation, rank 16 improves only modestly (\(\approx\)0.68\% relative FID gain from 10 to 20 epochs), without overturning the moderate-rank trend. Tiny DiT validation exhibits a similar ordering, with moderate ranks again outperforming rank 16. These results support a practical, reproducible guideline: under constrained and fixed budgets, small-to-moderate LoRA ranks are often a strong first choice for diffusion fine-tuning.
\end{abstract}
